%% Source: latex/notes.tex, sections 1-4 (Разрешимость, Спектр, Достижимые спектры, Критерий разрешимости).
%% Lean: NeoTiling.Core.{Neoclassical,Spectrum,Reachable,Solvability}.
\section{Solvability}\label{sec:solvability}

Throughout, $d\in\N$, $d\ge2$. For $p,x\in\R^d$ we write $p\circ x:=(p^1x^1,\dots,p^dx^d)$ for the coordinatewise product, $\e:=(1,\dots,1)\in\R^d$, and $\ind S\in\R^T$ for the indicator of $S\subset[T]$. Inequalities between vectors are coordinatewise. Prices $p_t\in\Rpp[d]$ are input prices divided by the output price of the same period, as in \cref{sec:intro}.

\begin{definition}[neoclassical functions]\label{def:neoclassical}
A function $f\colon\Rp[d]\to\R_+$ is \emph{neoclassical} if it is continuous, concave and positively homogeneous of degree one. The class of neoclassical functions of $d$ variables is denoted $\Phi_d$.
\end{definition}

A neoclassical function is nondecreasing on $\Rp[d]$ (see the proof of \cref{lem:necessity}). The class $\Phi_d$ is the class of unit cost functions of the generalized Houthakker--Johansen model, that is, of Young transforms of neoclassical production functions \citep{Shananin1997a}. It also contains functions that vanish on part of the orthant or have unbounded level sets, which \citet{Shananin1999} and \citet{Agaltsov2016} exclude; the criteria below are unaffected, since the function constructed in \cref{lem:patch} is positive on $\Rpp[d]$ with bounded level sets.

\begin{definition}[neoclassical solvability]\label{def:neo-solvability}
Let $\set{\hat p_t\in\Rpp[d]}_{t\in[T]}$ be a finite set of prices. An output vector $y\in\Rp[T]$ is \emph{neoclassically solvable} at these prices if there exist $h\in\Phi_d$ and neighborhoods $U_t\ni\hat p_t$, $U_t\subset\Rpp[d]$, such that for every choice $\set{p_t\in U_t}_{t\in[T]}$ there is a nonnegative absolutely continuous measure $\mu$ on $\Rp[d]$ with
\[
\mu(A_t)=y_t\quad\text{for all }t\in[T],\qquad A_t:=\set{x\in\Rpp[d]\mid h(p_t\circ x)<1}.
\]
\end{definition}

\subsection{Spectrum}

\begin{definition}[spectrum]\label{def:spectrum}
Let $h\in\Phi_d$ and $P:=\set{p_t\in\Rpp[d]}_{t\in[T]}$. For $S\subset[T]$ the \emph{chamber} with spectrum $S$ is
\[
R_S:=\set{x\in\Rpp[d]\mid h(p_t\circ x)<1\text{ for }t\in S,\ h(p_s\circ x)>1\text{ for }s\notin S}.
\]
The \emph{spectrum} of the pair $(h,P)$ is the family of spectra of nonempty chambers, $\Sp(h,P):=\set{S\subset[T]\mid R_S\ne\varnothing}$.
\end{definition}

The spectrum of a chamber lists the periods in which its technologies are profitable; this is the code of a cell in \citet{Shananin1999}.

\begin{proposition}\label{prop:null}
Let $h\in\Phi_d$. For every $p\in\Rpp[d]$ the set $\set{x\in\Rpp[d]\mid h(p\circ x)=1}$ has Lebesgue measure zero. Hence $\mathcal L^d\bigl(A_t\mathbin{\triangle}\bigsqcup_{S\ni t}R_S\bigr)=0$ for every $t\in[T]$.
\end{proposition}

\begin{proof}
The set $K:=\set{x\in\Rp[d]\mid h(p\circ x)\ge1}$ is convex since $h$ is concave. If $h(p\circ x)=1$, then $h(p\circ\lambda x)=\lambda<1$ for $0<\lambda<1$ by homogeneity, so $\lambda x\notin K$ and $\lambda x\to x$ as $\lambda\to1$. Hence the level set is contained in the boundary of the convex set $K$, and the boundary of a convex set has Lebesgue measure zero \citep{Lang1986}. The symmetric difference $A_t\mathbin{\triangle}\bigsqcup_{S\ni t}R_S$ is contained in the union of the level sets $\set{x\mid h(p_s\circ x)=1}$, $s\in[T]$.
\end{proof}

The following lemma is the solvability criterion of the moment problem for a fixed unit cost function \citep[Prop.~3.1]{Shananin1999}; we include the short proof because the measure it produces is used again in \cref{thm:solvability}.

\begin{lemma}[signature]\label{lem:signature}
Let $h\in\Phi_d$ and $P:=\set{p_t\in\Rpp[d]}_{t\in[T]}$. Then $y\in\Rp[T]$ lies in $\cone\set{\ind S\mid S\in\Sp(h,P)}$ if and only if there is an absolutely continuous measure $\mu$ on $\Rp[d]$ with $\mu(A_t)=y_t$ for all $t\in[T]$, where $A_t:=\set{x\in\Rpp[d]\mid h(p_t\circ x)<1}$.
\end{lemma}

\begin{proof}
By \cref{prop:null} the level sets are null, so $\mu(A_t)=\sum_{S\ni t}\mu(R_S)$, and the conditions $\mu(A_t)=y_t$ read $y=\sum_S\mu(R_S)\ind S$ with $\mu(R_S)=0$ for $S\notin\Sp(h,P)$. Conversely, for $S\in\Sp(h,P)$ the chamber $R_S$ is a nonempty open set, the uniform probability measure on a ball inside $R_S$ realizes $\ind S$, and sums and nonnegative multiples of measures realize sums and multiples of vectors.
\end{proof}

\begin{remark}[the sampled Radon transform]\label{rem:radon}
In the notation of \cref{sec:intro}, $\mu(A_t)=\partial_{p_0}\Pi(1,p_t)$ is the value at $p_t$ of the generalized Radon transform of $\mu$ \citep{Agaltsov2016}, and \cref{lem:signature} says that its values at finitely many points form the cone generated by the sampled indicators $\ind S$. With complete data the same cone is described by complete monotonicity after division by the Mellin symbol of $e^{-h}$ \citep{Agaltsov2015}; the measure constructed in the proof above, a finite sum of uniform measures on balls, shows that nothing of that analytic structure survives finite sampling: any solvable $y$ is realized by a compactly supported measure with a bounded density.
\end{remark}

\begin{lemma}[stability]\label{lem:stability}
Let $h\in\Phi_d$ and $P\in\Rpp[d\times T]$. There is a neighborhood $V\subset\Rpp[d\times T]$ of $P$ such that $\Sp(h,Q)\supseteq\Sp(h,P)$ for every $Q\in V$.
\end{lemma}

\begin{proof}
Let $S\in\Sp(h,P)$ and $x\in R_S$. The strict inequalities $h(p_t\circ x)\gtrless1$ survive a small change of prices by continuity of $h$, so $x$ lies in the chamber with spectrum $S$ for $Q$ as well, and $S\in\Sp(h,Q)$. There are finitely many $S$, so one neighborhood serves all of them.
\end{proof}

\subsection{Reachable spectra}

\begin{definition}[reachable spectra]\label{def:reachable}
Let $P:=\set{p_t\in\Rpp[d]}_{t\in[T]}$. A set $S\subset[T]$ is a \emph{reachable spectrum} at $P$ if for every $t\in S$ there is a nonzero $\xi\in\Rp[d]$ with $\ip{\xi}{p_s-p_t}>0$ for all $s\notin S$. Such $\xi$ is a \emph{witness} of the index $t$. The family of reachable spectra is denoted $\Sp(P)$.
\end{definition}

\begin{lemma}[necessity]\label{lem:necessity}
Let $h\in\Phi_d$ and $P:=\set{p_t\in\Rpp[d]}_{t\in[T]}$. Then $\Sp(h,P)\subseteq\Sp(P)$.
\end{lemma}

\begin{proof}
Let $S\in\Sp(h,P)$, $x\in R_S$ and $t\in S$. If $S=[T]$, any nonzero $\xi\in\Rp[d]$ is a witness. Otherwise the nonempty set $Z_x:=\conv\set{p_r\circ x}_{r\notin S}+\Rp[d]$ is convex and closed, and $h(z)\ge\min_{r\notin S}h(p_r\circ x)\ge1$ for all $z\in Z_x$ by concavity, homogeneity and monotonicity of $h$. Since $t\in S$ and $x\in R_S$ we have $h(p_t\circ x)<1$, so $p_t\circ x\notin Z_x$. Here $h$ is nondecreasing on $\Rp[d]$: by concavity and homogeneity it is superadditive, $h(z+v)=2h\bigl(\tfrac12(z+v)\bigr)\ge h(z)+h(v)\ge h(z)$ for $z,v\in\Rp[d]$, as $h\ge0$.

By the strong separation theorem for a point and a closed convex set \citep[Thm.~11.4]{Rockafellar1970} there are a linear functional $f$ and a number $u$ with $f(p_t\circ x)<u<f(z)$ for all $z\in Z_x$. The functional $f$ is nonnegative on $\Rp[d]$: if $f(v)<0$ for some $v\ge0$, then $z+\lambda v\in Z_x$ with a suitable $\lambda\ge0$ would give $f(z+\lambda v)=u$. Put $\xi^i:=f(e_i)x^i\ge0$. Then $\ip{\xi}{p}=f(p\circ x)$, and for $s\notin S$ we get $\ip{\xi}{p_s-p_t}=f(p_s\circ x)-f(p_t\circ x)>0$, whence $\xi\ne0$. Thus $\xi$ is a witness of $t$ at $P$, and $S\in\Sp(P)$.
\end{proof}

\subsection{The solvability criterion}

\begin{lemma}[separation]\label{lem:separation}
Let $S\in\Sp(P)$ and $t\in S$. There is $\xi\in\Rpp[d]$ with $\ip{\xi}{p_t}<1<\ip{\xi}{p_s}$ for all $s\notin S$.
\end{lemma}

\begin{proof}
Let $\xi_0$ be a witness of $t$. For small $\delta>0$ the vector $\xi_1:=\xi_0+\delta\e\in\Rpp[d]$ still satisfies $\ip{\xi_1}{p_s-p_t}>0$ for all $s\notin S$, there being finitely many such $s$. Choose $M$ with $\ip{\xi_1}{p_t}<M<\min_{s\notin S}\ip{\xi_1}{p_s}$, and $M:=\ip{\xi_1}{p_t}+1$ if $S=[T]$. Then $\xi:=\xi_1/M$ satisfies $\ip{\xi}{p_t}<1<\min_{s\notin S}\ip{\xi}{p_s}$.
\end{proof}

\begin{lemma}[patching]\label{lem:patch}
Let $P:=\set{p_t\in\Rpp[d]}_{t\in[T]}$. There is $h\in\Phi_d$ with $\Sp(P)\subseteq\Sp(h,P)$.
\end{lemma}

\begin{proof}
For every pair $(S,t)$ with $t\in S\in\Sp(P)$ take, by \cref{lem:separation}, $\xi_{S,t}\in\Rpp[d]$ with $\ip{\xi_{S,t}}{p_t}<1<\ip{\xi_{S,t}}{p_s}$ for $s\notin S$. Choose $M\ge1$ so large that
\[
M\min\bigl(\xi^1_{S,t}p^1_r,\ \xi^2_{S,t}p^2_r\bigr)>1\qquad\text{for all }S\in\Sp(P),\ t\in S,\ r\in[T].
\]
For $z\in\Z$ put $x_z:=(M^z,M^{-z},1,\dots,1)$. For $z\ne0$ one of $M^z$, $M^{-z}$ is at least $M$ and both terms are positive, so
\[
\ip{\xi_{S,t}}{p_r\circ x_z}\ge M\min\bigl(\xi^1_{S,t}p^1_r,\ \xi^2_{S,t}p^2_r\bigr)>1 .
\]
Fix an injection $E\colon\Sp(P)\to\N_+$ and put $\tilde\xi_{S,t}:=\xi_{S,t}\circ x_{-E_S}$ and
\[
h(x):=\min_{t\in S\in\Sp(P)}\ip{\tilde\xi_{S,t}}{x}\in\Phi_d .
\]
Since $\ip{\xi_{S,t}\circ x_{-E_S}}{p_r\circ x_{E_{S'}}}=\ip{\xi_{S,t}}{p_r\circ x_{E_{S'}-E_S}}$, for $S'\ne S$ every term in the minimum defining $h(p_r\circ x_{E_{S'}})$ with $S\ne S'$ exceeds $1$. Fix $S\in\Sp(P)$. For $r\notin S$ and every $t\in S$ we have $\ip{\xi_{S,t}}{p_r}>1$, hence $h(p_r\circ x_{E_S})>1$. For $t\in S$ we have $\ip{\xi_{S,t}}{p_t}<1$, hence $h(p_t\circ x_{E_S})<1$. Therefore $x_{E_S}\in R_S$ and $S\in\Sp(h,P)$.
\end{proof}

\begin{theorem}[solvability criterion]\label{thm:solvability}
Let $P:=\set{p_t\in\Rpp[d]}_{t\in[T]}$. An output vector $y\in\Rp[T]$ is neoclassically solvable at $P$ if and only if $y\in\cone\set{\ind S\mid S\in\Sp(P)}$.
\end{theorem}

\begin{proof}
Let $y$ be solvable with $h\in\Phi_d$. Taking the prices $P$ themselves in \cref{def:neo-solvability}, \cref{lem:signature} gives $y\in\cone\set{\ind S\mid S\in\Sp(h,P)}$, and $\Sp(h,P)\subseteq\Sp(P)$ by \cref{lem:necessity}. Conversely, take $h$ from \cref{lem:patch}. By \cref{lem:stability}, $\Sp(h,Q)\supseteq\Sp(h,P)\supseteq\Sp(P)$ for all $Q$ in a neighborhood of $P$, and by \cref{lem:signature} for every such $Q$ there is an absolutely continuous $\mu$ with $\mu(A_t)=y_t$.
\end{proof}

\begin{remark}[one input]\label{rem:d1}
For $d=1$ every neoclassical function is $h(x)=ax$ with $a\ge0$, so the sets $A_t$ are the intervals $(0,1/(ap_t))$, nested by price, and nothing depends on $a$. Accordingly $S\in\Sp(P)$ if and only if $p_s>p_t$ for all $t\in S$, $s\notin S$, and \cref{thm:solvability} reduces to: $y$ is solvable if and only if $y_s\ge y_t$ whenever $p_s<p_t$ and $y_s=y_t$ whenever $p_s=p_t$; weak solvability (\cref{sec:weak}) drops the second condition. The case $d=1$ is excluded from everything that follows.
\end{remark}

\begin{corollary}[openness]\label{lem:open}
Let $P\in\Rpp[d\times T]$. There is a neighborhood $V\subset\Rpp[d\times T]$ of $P$ such that $\Sp(Q)\supseteq\Sp(P)$ for every $Q\in V$.
\end{corollary}

\begin{proof}
Take $h\in\Phi_d$ with $\Sp(P)\subseteq\Sp(h,P)$ by \cref{lem:patch} and the neighborhood $V$ of \cref{lem:stability}. For $Q\in V$, \cref{lem:necessity} gives $\Sp(P)\subseteq\Sp(h,P)\subseteq\Sp(h,Q)\subseteq\Sp(Q)$.
\end{proof}

\begin{remark}[stability in place of uniqueness]
With complete data the profit function on an open cone of prices determines the distribution of capacities \citep[Thm.~2.1]{AgaltsovMolchanovShananin2018}, and the distribution is unique under a condition on the zeros of $\mathcal M[e^{-h}]$ \citep{Agaltsov2016}. With finite data neither $\mu$ nor the masses $\mu(R_S)$ are determined by $y$. \Cref{lem:open} is what remains: the set of prices at which a given $y$ is solvable is open, so the conclusion of \cref{thm:solvability} is insensitive to small errors in the prices, and \cref{def:neo-solvability} builds this in. Prices on the boundary of the openness are the subject of \cref{sec:weak}, and the distance of the data from this set is the subject of \cref{sec:lp}.
\end{remark}

